\documentclass[10pt,a4paper]{amsart}
\usepackage{amsmath,amssymb,amsthm,graphicx}
\usepackage{verbatim}
\newenvironment{reference}{\comment}{\endcomment}
%\newenvironment{reference}{}{}
\newenvironment{slogan}{\comment}{\endcomment}
\newenvironment{history}{\comment}{\endcomment}

% For commutative diagrams we use Xy-pic
\usepackage[all]{xy}

% We use 2cell for 2-commutative diagrams.
\xyoption{2cell}
% Theorem environments.
%
\theoremstyle{plain}
\newtheorem{theorem}[subsection]{Theorem}
\newtheorem{proposition}[subsection]{Proposition}
\newtheorem{lemma}[subsection]{Lemma}

\theoremstyle{definition}
\newtheorem{definition}[subsection]{Definition}
\newtheorem{example}[subsection]{Example}
\newtheorem{exercise}[subsection]{Exercise}
\newtheorem{situation}[subsection]{Situation}

\theoremstyle{remark}
\newtheorem{remark}[subsection]{Remark}
\newtheorem{remarks}[subsection]{Remarks}

\numberwithin{equation}{subsection}

% Macros
%
\def\lim{\mathop{\mathrm{lim}}\nolimits}
\def\colim{\mathop{\mathrm{colim}}\nolimits}
\def\Spec{\mathop{\mathrm{Spec}}}
\def\Hom{\mathop{\mathrm{Hom}}\nolimits}
\def\Ext{\mathop{\mathrm{Ext}}\nolimits}
\def\SheafHom{\mathop{\mathcal{H}\!\mathit{om}}\nolimits}
\def\SheafExt{\mathop{\mathcal{E}\!\mathit{xt}}\nolimits}
\def\Sch{\mathit{Sch}}
\def\Mor{\mathop{\mathrm{Mor}}\nolimits}
\def\Ob{\mathop{\mathrm{Ob}}\nolimits}
\def\Sh{\mathop{\mathit{Sh}}\nolimits}
\def\NL{\mathop{N\!L}\nolimits}
\def\CH{\mathop{\mathrm{CH}}\nolimits}
\def\proetale{{pro\text{-}\acute{e}tale}}
\def\etale{{\acute{e}tale}}
\def\QCoh{\mathit{QCoh}}
\def\Ker{\mathop{\mathrm{Ker}}}
\def\Im{\mathop{\mathrm{Im}}}
\def\Coker{\mathop{\mathrm{Coker}}}
\def\Coim{\mathop{\mathrm{Coim}}}

% Boxtimes
%
\DeclareMathSymbol{\boxtimes}{\mathbin}{AMSa}{"02}

%
% Macros for moduli stacks/spaces
%
\def\QCohstack{\mathcal{QC}\!\mathit{oh}}
\def\Cohstack{\mathcal{C}\!\mathit{oh}}
\def\Spacesstack{\mathcal{S}\!\mathit{paces}}
\def\Quotfunctor{\mathrm{Quot}}
\def\Hilbfunctor{\mathrm{Hilb}}
\def\Curvesstack{\mathcal{C}\!\mathit{urves}}
\def\Polarizedstack{\mathcal{P}\!\mathit{olarized}}
\def\Complexesstack{\mathcal{C}\!\mathit{omplexes}}
% \Pic is the operator that assigns to X its picard group, usage \Pic(X)
% \Picardstack_{X/B} denotes the Picard stack of X over B
% \Picardfunctor_{X/B} denotes the Picard functor of X over B
\def\Pic{\mathop{\mathrm{Pic}}\nolimits}
\def\Picardstack{\mathcal{P}\!\mathit{ic}}
\def\Picardfunctor{\mathrm{Pic}}
\def\Deformationcategory{\mathcal{D}\!\mathit{ef}}
\usepackage{cleveref}
\crefname{originallemma}{original lemma}{original lemmas}
\Crefname{originallemma}{Original lemma}{Original lemmas}
\makeatletter
\newlabel{algebra-lemma-NAK}{{00DV}{source}{}{sourcepointer-1}{}}
\newlabel{algebra-lemma-NAK@cref}{{[originallemma][0][]00DV}{[1][0][]source}{}{}{}}
\makeatother
\title{Surjectivity of a map of Artinian algebras is detected on the cotangent spaces}
\author{The Stacks Project Authors}
\date{stacks-project:formal-defos.tex:local-snapshot-20261012; GNU FDL 1.2 or later}
\begin{document}
\maketitle
\noindent\textit{Editorial note.} All delivered mathematical text is the author's own, copied from the pinned Stacks Project chapter file formal-defos.tex; no statement or proof step was rewritten, reordered or removed.\par\medskip
\noindent\textit{Editorial note.} The counted proof is the author's complete argument: the first implication through the nilpotence of the maximal ideals and the isomorphism on residue fields, the comparison of the absolute and relative cotangent spaces through the exact diagram, and the double application of Nakayama's lemma recovering surjectivity.\par\medskip
\noindent\textit{Editorial note.} The label lemma-surjective names the counted statement, and the proof environment immediately following it in the source is delivered verbatim as its proof.\par\medskip
\noindent\textit{Editorial note.} The counted statement is carried by the author's own lemma environment in the chapter source.\par\medskip
\noindent\textit{Editorial note.} Pre-open gate: python operations/agent-authoring/tools/precard\_gate.py --source formal-defos.tex --label lemma-surjective --start 563 --end 612 --local-id 28024 returned RESULT=OK (L28024 free, occupancy FREE, public interval clear). Candidate taken from sources/stacks-formal-defos/candidates-halfA.json.\par\medskip
\noindent\textit{Editorial note.} The statement and proof are copied from the pinned chapter file; the cited inputs are kept as external pointers with their canonical Stacks tags.\par\medskip
\noindent\textit{Allowed standard prerequisites (editorial).}\begin{itemize}
\item Nakayama's lemma. Algebra, the lemma allowing the surjectivity of a map of finite modules to be checked modulo a nilpotent ideal. Supplies the two reductions by which surjectivity of A -> B is checked on the quotients by the maximal ideal of A.
\item Exactness of the cotangent sequence for the residue field. The standard exact sequence relating the cotangent space of Lambda, that of an Artinian algebra and the relative cotangent space. Supplies the displayed commutative diagram with exact rows, which transfers surjectivity between the absolute and relative cotangent spaces.
\end{itemize}\medskip
\ifdefined\hypertarget\hypertarget{sourcepointer-1}{}\fi
\noindent\textit{Source reference.} Original lemma 00DV: sources/stacks-chow/tags.txt, source lines 512--512. Prerequisite pointer invoked at proof line(s) 604; it names a standard fact of a sibling chapter that the counted argument uses, and every step of the counted proof built on it is retained verbatim in the delivered body.\par\medskip
\begin{lemma}
\label{lemma-surjective}
Let $A \to B$ be a ring map in $\mathcal{C}_\Lambda$.
The following are equivalent
\begin{enumerate}
\item $f$ is surjective,
\item $\mathfrak m_A/\mathfrak m_A^2 \to \mathfrak m_B/\mathfrak m_B^2$
is surjective, and
\item $\mathfrak m_A/(\mathfrak m_\Lambda A + \mathfrak m_A^2)
\to \mathfrak m_B/(\mathfrak m_\Lambda B + \mathfrak m_B^2)$ is surjective.
\end{enumerate}
\end{lemma}

\begin{proof}
For any ring map $f : A \to B$ in $\mathcal{C}_\Lambda$ we have
$f(\mathfrak m_A) \subset \mathfrak m_B$ for example because
$\mathfrak m_A$, $\mathfrak m_B$ is the set of nilpotent elements of
$A$, $B$. Suppose $f$ is surjective. Let $y \in \mathfrak m_B$.
Choose $x \in A$ with $f(x) = y$. Since $f$ induces an isomorphism
$A/\mathfrak m_A \to B/\mathfrak m_B$ we see that $x \in \mathfrak m_A$.
Hence the induced map
$\mathfrak m_A/\mathfrak m_A^2 \to \mathfrak m_B/\mathfrak m_B^2$
is surjective. In this way we see that (1) implies (2).

\medskip\noindent
It is clear that (2) implies (3). The map $A \to B$ gives rise
to a canonical commutative diagram
$$
\xymatrix{
\mathfrak m_\Lambda/\mathfrak m_\Lambda^2 \otimes_{k'} k \ar[r] \ar[d] &
\mathfrak m_A/\mathfrak m_A^2 \ar[r] \ar[d] &
\mathfrak m_A/(\mathfrak m_\Lambda A + \mathfrak m_A^2) \ar[r] \ar[d] & 0 \\
\mathfrak m_\Lambda/\mathfrak m_\Lambda^2 \otimes_{k'} k \ar[r] &
\mathfrak m_B/\mathfrak m_B^2 \ar[r] &
\mathfrak m_B/(\mathfrak m_\Lambda B + \mathfrak m_B^2) \ar[r] & 0
}
$$
with exact rows. Hence if (3) holds, then so does (2).

\medskip\noindent
Assume (2). To show that $A \to B$ is surjective it suffices by
Nakayama's lemma (Algebra, Lemma \ref{algebra-lemma-NAK})
to show that $A/\mathfrak m_A \to B/\mathfrak m_AB$ is surjective.
(Note that $\mathfrak m_A$ is a nilpotent ideal.)
As $k = A/\mathfrak m_A = B/\mathfrak m_B$ it suffices to show that
$\mathfrak m_AB \to \mathfrak m_B$ is surjective. Applying
Nakayama's lemma once more we see that it suffices to see that
$\mathfrak m_AB/\mathfrak m_A\mathfrak m_B \to \mathfrak m_B/\mathfrak m_B^2$
is surjective which is what we assumed.
\end{proof}

\end{document}
